\chapter{The Rates Options Market}\label{ch:m2:the-rates-options-market}

On 11 June 2014 the European Central Bank cut the rate it pays on deposits to
$-0.10\%$, and over the next five years to $-0.50\%$. German ten-year yields
were below zero in 38 months between mid-2016 and early 2022. For the options
market this was not a curiosity but a breakdown. The model long standard for
these options, Black's, assumes that a
rate moves in proportion to its level, and a rate of $0.10\%$ that rises to
$0.20\%$ has moved 100\%; below zero the model has no answer at all. The market
switched to quoting in the units a rates trader thinks in, basis points of
movement, with a model written in 1900 by a French doctoral student. This
chapter explains the options traded on interest rates, \omterm{def:m2:the-rates-options-market:cap}{caps}, \omterm{def:m2:the-rates-options-market:cap}{floors} and
\omterm{def:m2:the-rates-options-market:swaption}{swaptions}, the \omterm{def:m2:the-rates-options-market:normal}{normal volatility} in which they are quoted, and the structural
flows, from \omterm{def:m2:the-rates-options-market:swaption}{callable bonds} and mortgages, that set their prices.

\section{Caps and floors}

\begin{definition}[Caplet, cap, floor]\label{def:m2:the-rates-options-market:cap}
A \emph{caplet}\index{caplet} with strike $K$ on a floating rate for a period of
length $\delta$ pays $N\delta\max(L - K, 0)$ at the end of the period, where $L$
is the rate fixed for that period. A \emph{cap}\index{cap} is a strip of
caplets on consecutive periods; a \emph{floor}\index{floor} is the same strip of
floorlets, each paying $N\delta\max(K - L, 0)$.
\end{definition}

A borrower who pays a floating rate buys a \omterm{def:m2:the-rates-options-market:cap}{cap} to limit what it can pay; a
lender, or an investor in floating-rate notes, buys a \omterm{def:m2:the-rates-options-market:cap}{floor}. Buying a \omterm{def:m2:the-rates-options-market:cap}{cap} and
selling a \omterm{def:m2:the-rates-options-market:cap}{floor} is a collar, often struck so that the two premiums cancel. Each
\omterm{def:m2:the-rates-options-market:cap}{caplet} is an option on one forward rate, so a \omterm{def:m2:the-rates-options-market:cap}{cap} is priced as a sum of
single-period options with one volatility per period, or one flat volatility
quoted for the whole \omterm{def:m2:the-rates-options-market:cap}{cap}.

\begin{example}[A cap and a zero-cost collar]\label{ex:m2:the-rates-options-market:cap}
On a flat 4\% curve, a five-year \omterm{def:m2:the-rates-options-market:cap}{cap} at 4.5\% on USD~100 million of an annual
rate, with the first period already fixed, has four \omterm{def:m2:the-rates-options-market:cap}{caplets} fixing in one to
four years. At a \omterm{def:m2:the-rates-options-market:normal}{normal volatility} of 100 basis points a year they are worth
USD~182\,874, 310\,339, 401\,400 and 470\,709: USD~1.37 million in all, 137
basis points of notional. A \omterm{def:m2:the-rates-options-market:cap}{floor} at 3.5\% is worth exactly the same, since the
forward sits halfway between the strikes and the normal model is symmetric: the
borrower who buys the \omterm{def:m2:the-rates-options-market:cap}{cap} and sells the \omterm{def:m2:the-rates-options-market:cap}{floor} pays nothing and keeps its rate
between 3.5\% and 4.5\%.
\end{example}

\section{Swaptions}

\begin{definition}[Swaption, payer, receiver; callable bond]\label{def:m2:the-rates-options-market:swaption}
A \emph{swaption}\index{swaption} is an option to enter an \omterm{def:m2:interest-rate-swaps:swap}{interest-rate swap}
at a fixed rate $K$ on a future date. A \emph{payer
swaption}\index{payer swaption} gives the right to pay fixed; a
\emph{receiver swaption}\index{receiver swaption} the right to receive fixed.
At expiry $T$ a payer is worth $N A(T)\max(S_T - K, 0)$, where $S_T$ is the par
rate of the underlying swap then and $A(T)$ its annuity. A \emph{callable
bond}\index{callable bond} is a bond its issuer may repay early, at a set price
on set dates.
\end{definition}

\omterm{def:m2:the-rates-options-market:swaption}{Swaptions} are named by expiry and tenor: a 1y$\times$10y is a one-year option on
a ten-year swap. A payer is a bet on higher rates, a receiver on lower rates;
the pair struck at the forward, the straddle, is a bet on how much rates move,
in either direction (\cref{fig:m2:the-rates-options-market:payoff}). The
annuity turns a rate difference into money: pricing uses the forward swap rate
and the forward annuity, and in that measure the swap rate is a martingale, so
the \omterm{def:m2:the-rates-options-market:swaption}{swaption} is priced like an option on a single rate.

\begin{omfigure}
\begin{tikzpicture}
\begin{axis}[width=11cm,height=4.8cm,
  xlabel={ten-year swap rate at expiry (\%)},ylabel={value per 100 notional},
  tick label style={font=\small},label style={font=\small},
  xmin=2,xmax=6,ymin=0,ymax=16,grid=major,grid style={gray!25},
  legend columns=3,legend style={font=\footnotesize,at={(0.5,-0.32)},anchor=north,draw=none,fill=none,/tikz/every even column/.append style={column sep=8pt}},
  table/col sep=comma]
\addplot[omThm,very thick] table[x=s,y=payer]{figdata/markets-2/13-the-rates-options-market/payoff.csv};
\addplot[omDef,very thick] table[x=s,y=receiver]{figdata/markets-2/13-the-rates-options-market/payoff.csv};
\addplot[omThm,thick,dashed] table[x=s,y=payer_value]{figdata/markets-2/13-the-rates-options-market/payoff.csv};
\legend{payer at expiry,receiver at expiry,payer a year before expiry}
\end{axis}
\end{tikzpicture}

\omcaption{A 1y$\times$10y payer and \omterm{def:m2:the-rates-options-market:swaption}{receiver swaption} struck at a 4\% forward,
per 100 of notional, on a flat 4\% curve. At expiry each pays the annuity (7.80)
times the distance of the swap rate beyond the strike; a year before, the payer
is worth its time value on top, at a \omterm{def:m2:the-rates-options-market:normal}{normal volatility} of 95 basis points.
Illustrative; data: the chapter's tutorial.}\label{fig:m2:the-rates-options-market:payoff}
\end{omfigure}

\omterm{def:m2:the-rates-options-market:swaption}{Callable bonds} are one natural source of \omterm{def:m2:the-rates-options-market:swaption}{swaption} supply. An issuer that
sells a ten-year bond callable in three years holds the right to stop paying its
coupon from year three, which is a \omterm{def:m2:the-rates-options-market:swaption}{receiver swaption}: when rates fall, it calls
the bond and refinances. Issuers that fund at a floating rate swap the bond
with a dealer and sell that right to the dealer inside the swap, which the
dealer can cancel on the call date; the dealer is then long a \omterm{def:m2:the-rates-options-market:swaption}{receiver
swaption}, and long volatility, and sells \omterm{def:m2:the-rates-options-market:swaption}{swaptions} to other investors to lay
the risk off (\cref{fig:m2:the-rates-options-market:callable}). The weekend
problem does this trade.

\begin{omfigure}
\centering
\begin{tikzpicture}[font=\small,
  box/.style={draw,thick,rounded corners=2pt,align=center,minimum height=1cm,minimum width=2.4cm,inner sep=3pt},
  arr/.style={-{Stealth[length=2.5mm]},thick}]
\node[box,draw=omDef,fill=omDef!8] (inv) at (0,0) {bond investors};
\node[box,draw=omThm,fill=omThm!8] (iss) at (4.6,0) {issuer};
\node[box,draw=omThm,fill=omThm!8] (dl) at (9.2,0) {dealer};
\node[box,draw=omProp,fill=omProp!8] (vb) at (13.2,0) {volatility\\buyers};
\draw[arr] (inv) -- node[above=5pt,font=\footnotesize,align=center] {buy a\\callable at a\\higher coupon} (iss);
\draw[arr] (iss) -- node[above=5pt,font=\footnotesize,align=center] {cancellable\\swap: sells\\the call} (dl);
\draw[arr] (dl) -- node[above=5pt,font=\footnotesize,align=center] {sells\\swaptions} (vb);
\node[font=\footnotesize,text=omThm,align=center] at (4.6,-1.2) {pays floating};
\node[font=\footnotesize,text=omThm,align=center] at (9.2,-1.25) {long a receiver swaption,\\long vega until laid off};
\end{tikzpicture}

\omcaption{How the option in a \omterm{def:m2:the-rates-options-market:swaption}{callable bond} reaches the volatility market.
Investors are paid for the issuer's call in the coupon; the issuer sells the
same option to a dealer inside a cancellable swap and pays floating; the dealer
sells \omterm{def:m2:the-rates-options-market:swaption}{swaptions} to investors who want volatility. Arrows show who sells the
option to whom. Schematic.}\label{fig:m2:the-rates-options-market:callable}
\end{omfigure}

Mortgages run the other way. The holder of a mortgage-backed security is short
the borrowers' refinancing option (\cref{ch:m2:mortgages-and-agencies}), and
some hedge that negative convexity by buying \omterm{def:m2:the-rates-options-market:swaption}{swaptions}. Issuers of callable
debt are natural sellers of volatility, mortgage investors natural buyers, and
dealers stand between them.

\section{Normal volatility}

\begin{definition}[Bachelier model, normal volatility]\label{def:m2:the-rates-options-market:normal}
The \emph{Bachelier model}\index{Bachelier model}, published by Louis Bachelier
in 1900, assumes that the forward rate follows $dF_t = \sigma_N\,dW_t$: it
moves by normally distributed amounts, independent of its level, and can go
negative. Its parameter $\sigma_N$, in basis points a year, is the \emph{normal
volatility}\index{normal volatility}; the implied normal volatility of an option
is the $\sigma_N$ at which the model reproduces its price.
\end{definition}

\begin{proposition}[The Bachelier price]\label{prop:m2:the-rates-options-market:bachelier}
In the \omterm{def:m2:the-rates-options-market:normal}{Bachelier model} a payer with strike $K$ and expiry $T$, on a forward
rate $F$ with annuity $A$, is worth
\[
A\left[(F-K)\,\Phi(d) + \sigma_N\sqrt{T}\,\varphi(d)\right], \qquad d = \frac{F-K}{\sigma_N\sqrt T},
\]
and a receiver is the payer less $A(F-K)$. At the money each is worth
$A\sigma_N\sqrt{T}/\sqrt{2\pi}$, so a straddle costs
$A\sigma_N\sqrt{T}\sqrt{2/\pi}$.
\end{proposition}

\begin{proof}
$F_T$ is normal with mean $F$ and standard deviation $s = \sigma_N\sqrt{T}$.
With $F_T = F + sZ$, $\mathbb{E}[(F_T-K)^+] = \mathbb{E}[(F-K+sZ)\mathbf 1\{Z>-d\}]
= (F-K)\Phi(d) + s\varphi(d)$, using $\mathbb{E}[Z\mathbf 1\{Z>-d\}] = \varphi(d)$.
Put--call parity gives the receiver, and $d = 0$ the at-the-money value.
\end{proof}

A trader converts the annual \omterm{def:m2:the-rates-options-market:normal}{normal volatility} into a daily one to compare it
with what rates actually do: 95 basis points a year is
$95/\sqrt{252} = 5.98$ basis points a business day. If the ten-year rate moves
more than six basis points a day on average over the option's life, a
delta-hedged straddle bought at 95 makes money; less, and it loses. Quoting in
\omterm{def:m2:the-rates-options-market:normal}{normal volatility} makes this comparison immediate, and it remains the usual
convention for rates options now that rates are positive again.

The lognormal, Black, quote is still used and must be converted. At the money,
$\sigma_N \approx \sigma_B F$; away from it the models' smiles differ. Near zero
the conversion explodes (\cref{fig:m2:the-rates-options-market:black}): 90
basis points of \omterm{def:m2:the-rates-options-market:normal}{normal volatility} is a Black volatility of 22.5\% at a 4\%
forward, 93\% at 1\% and 215\% at 0.5\%, and below a forward of about 0.36\%
no Black volatility at all reproduces the price, because a Black call can never
be worth more than the forward. When euro and yen rates spent years near and
below zero (\cref{fig:m2:the-rates-options-market:negative}), Black quotes
stopped carrying information and the normal quote took over.

\begin{omfigure}
\begin{tikzpicture}
\begin{axis}[width=11cm,height=4.8cm,ymode=log,
  xlabel={forward rate (\%)},ylabel={Black volatility (\%)},
  tick label style={font=\small},label style={font=\small},
  xmin=0,xmax=6.2,ymin=10,ymax=500,grid=major,grid style={gray!25},
  log ticks with fixed point,
  table/col sep=comma]
\addplot[omThm,very thick,mark=*,mark size=1.2pt] table[x=forward,y=black]{figdata/markets-2/13-the-rates-options-market/blackequiv.csv};
\draw[omProp,thick,dashed] (axis cs:0.359,10) -- (axis cs:0.359,500);
\node[omProp,font=\footnotesize,anchor=west,align=left] (nb) at (axis cs:0.9,19) {no Black volatility\\below 0.36\%};
\draw[omProp,thick,-{Stealth}] (nb.west) -- (axis cs:0.40,19);
\end{axis}
\end{tikzpicture}

\omcaption{The Black volatility that gives the same one-year at-the-money price
as a \omterm{def:m2:the-rates-options-market:normal}{normal volatility} of 90 basis points, by forward rate, on a logarithmic
scale. The same option looks calm at 5\% and wild at 0.5\%; below the dashed
line no Black volatility exists. Data: the chapter's
tutorial.}\label{fig:m2:the-rates-options-market:black}
\end{omfigure}

\begin{omfigure}
\begin{tikzpicture}
\begin{axis}[width=11cm,height=4.6cm,
  ylabel={10-year yield (\%)},
  tick label style={font=\small},label style={font=\small},
  xmin=2012,xmax=2023,ymin=-0.8,ymax=2.2,grid=major,grid style={gray!25},
  xtick={2012,2014,2016,2018,2020,2022},xticklabel style={/pgf/number format/1000 sep={}},
  legend columns=2,legend style={font=\footnotesize,at={(0.5,-0.2)},anchor=north,draw=none,fill=none,/tikz/every even column/.append style={column sep=8pt}},
  table/col sep=comma]
\addplot[omDef,very thick] table[x=t,y=DE]{figdata/markets-2/13-the-rates-options-market/negyields.csv};
\addplot[omThm,thick,dashed] table[x=t,y=JP]{figdata/markets-2/13-the-rates-options-market/negyields.csv};
\draw[gray] (axis cs:2012,0) -- (axis cs:2023,0);
\legend{Germany,Japan}
\end{axis}
\end{tikzpicture}

\omcaption{German and Japanese ten-year government yields, monthly averages,
2012 to 2022. The German yield was negative in 38 months between June 2016 and
January 2022, the Japanese in 24 between February 2016 and April 2020. Data:
OECD via FRED series IRLTLT01DEM156N and
IRLTLT01JPM156N.}\label{fig:m2:the-rates-options-market:negative}
\end{omfigure}

\begin{example}[A straddle]\label{ex:m2:the-rates-options-market:straddle}
On a flat 4\% curve the 1y$\times$10y forward annuity is 7.80. At 95 basis
points of \omterm{def:m2:the-rates-options-market:normal}{normal volatility}, a straddle on USD~100 million costs
$7.80 \times 0.0095 \times \sqrt{2/\pi} = 5.91\%$ of notional, USD~5.91 million.
Divided by the annuity it is 75.8 basis points: the ten-year rate must end
more than about 76 basis points away from 4\% for the straddle held to expiry to
pay for itself.
\end{example}

\section{The volatility cube and who trades it}

\begin{definition}[Volatility cube]\label{def:m2:the-rates-options-market:cube}
The \emph{volatility cube}\index{volatility cube} of a currency is the set of
implied volatilities of its \omterm{def:m2:the-rates-options-market:swaption}{swaptions} by expiry, tenor of the underlying swap and
strike. Its at-the-money face, expiry by tenor, is the \omterm{def:m2:the-rates-options-market:swaption}{swaption} matrix; the
strike dimension is the smile.
\end{definition}

A rates options desk quotes and risk-manages the whole cube. Its at-the-money
matrix has a term structure in each direction
(\cref{fig:m2:the-rates-options-market:cube}): short expiries on short tenors
move with the next central-bank decisions, long expiries on long tenors with
the slow flows of callable issuance and mortgage hedging. The smile, which a
model such as SABR (One Quant Book 6) interpolates, prices the tails; for rates
it is usually higher on both sides of the money, and asymmetric.

\begin{omfigure}
\begin{tikzpicture}
\begin{axis}[name=left,width=6.4cm,height=4.8cm,xmode=log,
  xlabel={expiry (years)},ylabel={normal vol (bp a year)},
  tick label style={font=\small},label style={font=\small},
  xmin=0.2,xmax=12,ymin=70,ymax=120,grid=major,grid style={gray!25},
  xtick={0.25,1,5,10},xticklabels={3m,1y,5y,10y},
  legend columns=3,legend style={font=\footnotesize,at={(1.08,-0.3)},anchor=north,draw=none,fill=none,/tikz/every even column/.append style={column sep=8pt}},
  table/col sep=comma]
\addplot[omDef,very thick,mark=*,mark size=1.2pt] table[x=expiry,y=t2]{figdata/markets-2/13-the-rates-options-market/cube.csv};
\addplot[omThm,very thick,mark=*,mark size=1.2pt] table[x=expiry,y=t10]{figdata/markets-2/13-the-rates-options-market/cube.csv};
\addplot[omProp,very thick,mark=*,mark size=1.2pt] table[x=expiry,y=t30]{figdata/markets-2/13-the-rates-options-market/cube.csv};
\legend{2-year tenor,10-year tenor,30-year tenor}
\end{axis}
\begin{axis}[at={(left.east)},anchor=west,xshift=1.4cm,width=6.4cm,height=4.8cm,
  xlabel={strike minus forward (bp)},
  tick label style={font=\small},label style={font=\small},
  xmin=-210,xmax=210,ymin=90,ymax=112,grid=major,grid style={gray!25},
  table/col sep=comma]
\addplot[omThm,very thick,mark=*,mark size=1.2pt] table[x=offset,y=vol]{figdata/markets-2/13-the-rates-options-market/smile.csv};
\end{axis}
\end{tikzpicture}

\omcaption{Two faces of an illustrative dollar \omterm{def:m2:the-rates-options-market:cube}{volatility cube}: at-the-money
\omterm{def:m2:the-rates-options-market:normal}{normal volatility} by expiry for three tenors (left), and the smile of the
1y$\times$10y \omterm{def:m2:the-rates-options-market:swaption}{swaption} by strike (right). The shapes, not the numbers, are the
point. Data: the chapter's tutorial.}\label{fig:m2:the-rates-options-market:cube}
\end{omfigure}

The flows that shape the cube can be large and one-sided. In Taiwan, a law of
May 2014 let life insurers hold locally issued foreign-currency bonds outside
their \omterm{def:m2:the-rates-options-market:cap}{cap} on overseas assets; every dollar bond sold there in the following
months was callable by its issuer, with maturities of twenty or thirty years,
and three in four paid no coupon. Each such issue, swapped by its issuer, leaves
a dealer long a long-dated \omterm{def:m2:the-rates-options-market:swaption}{receiver swaption} that it must sell on: a steady
supply of long-dated volatility, from the balance sheets of one country's
insurers.

\begin{dated}{2026-09}{The market's size}\label{dat:m2:the-rates-options-market:size}
Over-the-counter derivatives notional outstanding was USD~846 trillion at the
end of June 2025, 16\% more than a year earlier, 79\% of it in interest-rate
derivatives; their gross market value was USD~21.8 trillion (BIS statistics
published 8 December 2025). \omterm{def:m2:the-rates-options-market:swaption}{Swaptions}, \omterm{def:m2:the-rates-options-market:cap}{caps} and \omterm{def:m2:the-rates-options-market:cap}{floors} are a part of the
interest-rate total; the BIS breaks options out in its detailed tables.
\end{dated}

\section{Tutorial: normal and lognormal quotes}\label{tut:m2:the-rates-options-market:quotes}

\begin{tutorial}
\textbf{Goal.} Price \omterm{def:m2:the-rates-options-market:cap}{caps} and \omterm{def:m2:the-rates-options-market:swaption}{swaptions} in the Bachelier and Black models,
convert quotes between them, find where the conversion fails, and price a
straddle and a \omterm{def:m2:the-rates-options-market:swaption}{callable bond}'s option from a cube slice. \textbf{End state:}
\cref{fig:m2:the-rates-options-market:payoff,fig:m2:the-rates-options-market:black},
\cref{ex:m2:the-rates-options-market:cap,ex:m2:the-rates-options-market:straddle}
and the numbers of the weekend problem.
\begin{tutsteps}
\item \textbf{The two pricers and the vega} of
\cref{prop:m2:the-rates-options-market:bachelier}.
\omcode{code/firm/normalvol/firm_normalvol.py}{20}{43}{Bachelier and Black prices, normal vega.}\label{lst:m2:the-rates-options-market:pricers}
\item \textbf{Implied volatility and conversion}, by bisection, refusing a
price that no volatility reproduces.
\omcode{code/firm/normalvol/firm_normalvol.py}{46}{68}{Implied volatility and quote conversion.}\label{lst:m2:the-rates-options-market:implied}
\item \textbf{Run} \texttt{swaption\_demo.straddle\_1y10y()},
\texttt{swaption\_demo.cap()}, \texttt{swaption\_demo.callable\_swap()} and
\texttt{fig\_swaptions.py}.
\end{tutsteps}
\textbf{What to change next.} Convert the smile of
\cref{fig:m2:the-rates-options-market:cube} into Black volatilities at a 4\%
and at a 1\% forward, and compare the shapes; then shift the curve down 3.5
points and repeat.
\end{tutorial}

\section{Build: the normal-volatility pricer}\label{bld:m2:the-rates-options-market:normalvol}

\begin{build}
\bfield{Purpose} The miniature firm quotes \omterm{def:m2:the-rates-options-market:cap}{caps}, \omterm{def:m2:the-rates-options-market:cap}{floors} and \omterm{def:m2:the-rates-options-market:swaption}{swaptions} in \omterm{def:m2:the-rates-options-market:normal}{normal
volatility}, receives Black quotes from some counterparties and must convert
them, and hedges vega from structured trades.
\bfield{Interface} \texttt{bachelier(f, k, t, vol, annuity, payer)};
\texttt{black(\ldots)}; \texttt{normal\_vega}; \texttt{implied(price, f, k, t,
annuity, payer, model)}; \texttt{normal\_to\_black};
\texttt{black\_to\_normal}; \texttt{atm\_straddle(t, vol, annuity)};
\texttt{bp\_per\_day(vol)}.
\bfield{Rules} Rates and volatilities as decimals; Black refuses a non-positive
forward or strike; implied volatility by bisection between explicit bounds,
raising an error when the price is outside the model's range.
\bfield{Acceptance tests} \texttt{code/firm/normalvol/tests/}: put--call parity in
both models, including negative rates; the at-the-money closed form; vega
against a finite difference; implied-volatility and conversion round trips;
the refusal at a 0.25\% forward.
\bfield{Stretch} Shifted Black; SABR in its normal form (One Quant Book 6);
Bermudan \omterm{def:m2:the-rates-options-market:swaption}{swaptions} by a lattice (One Quant Book 6); a cube interpolator by
expiry and tenor.
\end{build}

\begin{omsources}
\item L.~Bachelier, ``Th\'eorie de la sp\'eculation'', \emph{Annales
scientifiques de l'\'Ecole Normale Sup\'erieure}, 1900.
\item European Central Bank, key ECB interest rates; OECD long-term interest
rates via FRED.
\item F.~Le Floc'h, ``Explicit rational formulae for Bachelier (normal) implied
volatility'', arXiv 2605.18343, 2026.
\item Bloomberg, ``Goldman leads Wall Street tapping Taiwan cash pile'', 10
September 2014 (via Insurance Journal).
\item Bank for International Settlements, OTC derivatives statistics at
end-June 2025.
\end{omsources}

\section{Exercises}

\begin{exercise}[$\star$]\label{exo:m2:the-rates-options-market:1}
Express a \omterm{def:m2:the-rates-options-market:normal}{normal volatility} of 95 basis points a year in basis points a
business day, and a daily volatility of 7 basis points as an annual one.
\end{exercise}

\begin{exercise}[$\star$]\label{exo:m2:the-rates-options-market:2}
A company has a floating-rate loan. Which does it buy to protect itself, a payer
or a \omterm{def:m2:the-rates-options-market:swaption}{receiver swaption}, a \omterm{def:m2:the-rates-options-market:cap}{cap} or a \omterm{def:m2:the-rates-options-market:cap}{floor}? And a pension fund that fears lower
rates?
\end{exercise}

\begin{exercise}[$\star$]\label{exo:m2:the-rates-options-market:3}
Why can the Black model not price an option on a rate of $-0.2\%$, and why
did the market choose \omterm{def:m2:the-rates-options-market:normal}{normal volatility} rather than another fix?
\end{exercise}

\begin{exercise}[$\star\star$]\label{exo:m2:the-rates-options-market:4}
Give the price of the 1y$\times$10y at-the-money payer on USD~100 million at
95 basis points, and the straddle's break-even move.
\end{exercise}

\begin{exercise}[$\star\star$]\label{exo:m2:the-rates-options-market:5}
Give the Black volatilities equivalent to 90 basis points of \omterm{def:m2:the-rates-options-market:normal}{normal volatility},
one year at the money, at forwards of 6\%, 2\% and 0.4\%.
\end{exercise}

\begin{exercise}[$\star\star$]\label{exo:m2:the-rates-options-market:6}
In \cref{ex:m2:the-rates-options-market:cap}, why is the \omterm{def:m2:the-rates-options-market:cap}{floor} at 3.5\% worth
exactly the \omterm{def:m2:the-rates-options-market:cap}{cap} at 4.5\%? Would it still be so in Black's model?
\end{exercise}

\begin{exercise}[$\star\star\star$]\label{exo:m2:the-rates-options-market:7}
\emph{Coding.} Find the forward below which no Black volatility reproduces the
one-year at-the-money price at 90 basis points of \omterm{def:m2:the-rates-options-market:normal}{normal volatility}, and check
it against $\sigma_N\sqrt T/\sqrt{2\pi}$.
\end{exercise}

\begin{exercise}[$\star\star\star$]\label{exo:m2:the-rates-options-market:8}
\emph{Find the flaw.} ``\omterm{def:m2:the-rates-options-market:swaption}{Swaption} volatility is 20\% in one currency and 60\% in
another: the second is three times riskier.'' Correct it.
\end{exercise}

\section{Problem: The Callable Issuer}

\begin{problem}[{Weekend problem --- the vega a bank takes on from one callable bond}]\label{pb:m2:the-rates-options-market:1}
A bank issues USD~500 million of ten-year bonds with a 4.5\% annual coupon,
callable at par once, in three years. It swaps them with a dealer into
floating: the dealer pays 4.5\% fixed and receives floating for ten years, and
may cancel the swap in three years, when the bank would call the bond. The
curve is flat at 4\%; the 3y$\times$7y \omterm{def:m2:the-rates-options-market:normal}{normal volatility} is 100 basis points.

\medskip\noindent\textbf{Part I --- The option.}
\begin{enumerate}
\item What option does the bank hold through its bond, and why?
\item What option does the dealer hold through its swap?
\item Give the forward swap rate and the forward annuity of the 3y$\times$7y swap.
\item Is the option in or out of the money, and by how much?
\item Give its value on USD~500 million, in dollars and in basis points of notional.
\end{enumerate}

\medskip\noindent\textbf{Part II --- The dealer's risk.}
\begin{enumerate}[resume]
\item Give its vega per basis point of \omterm{def:m2:the-rates-options-market:normal}{normal volatility}.
\item Give its delta per basis point of the forward swap rate, and its sign.
\item How much does the dealer lose if volatility falls 2 basis points?
\item Give the vega of USD~100 million of at-the-money 3y$\times$7y straddles.
\item What notional of straddles must the dealer sell to lay off the vega?
\end{enumerate}

\medskip\noindent\textbf{Part III --- The bank.}
\begin{enumerate}[resume]
\item How does the bank get paid for the option it sells?
\item Express the option's value as a running spread over the ten years.
\item When will the bank call the bond, and what happens to the swap?
\item Why does the bank prefer this to a plain bond and a plain swap?
\item What happens to the dealer's hedge if rates fall to 3.5\% the next day?
\end{enumerate}

\medskip\noindent\textbf{Part IV --- The market.}
\begin{enumerate}[resume]
\item What happens to long-dated \omterm{def:m2:the-rates-options-market:swaption}{swaption} volatility when many issuers do this?
\item Who buys the volatility the dealers sell?
\item Why is a real \omterm{def:m2:the-rates-options-market:swaption}{callable bond}, callable every year, harder to hedge?
\item State the \emph{named result}: the vega the dealer acquires, and the
straddles it sells.
\item In one sentence: where does the supply of long-dated volatility come from?
\end{enumerate}
\end{problem}

\section{Interview questions}

\begin{interviewq}[$\star$ \iqroles{trader, researcher}]\label{iq:m2:the-rates-options-market:1}
What is the difference between normal and lognormal volatility, and when does
it matter?
\end{interviewq}

\begin{interviewq}[$\star$ \iqroles{trader, bank}]\label{iq:m2:the-rates-options-market:2}
Explain a payer and a \omterm{def:m2:the-rates-options-market:swaption}{receiver swaption}, and who buys each.
\end{interviewq}

\begin{interviewq}[$\star\star$ \iqroles{researcher}]\label{iq:m2:the-rates-options-market:3}
Derive the price of an at-the-money option in the \omterm{def:m2:the-rates-options-market:normal}{Bachelier model}.
\end{interviewq}

\begin{interviewq}[$\star\star$ \iqroles{trader}]\label{iq:m2:the-rates-options-market:4}
You are long a delta-hedged 1y$\times$10y straddle at 95 basis points. How do
you know, day by day, whether you are making money?
\end{interviewq}

\begin{interviewq}[$\star\star$ \iqroles{bank, trader}]\label{iq:m2:the-rates-options-market:5}
Why does callable issuance lower long-dated \omterm{def:m2:the-rates-options-market:swaption}{swaption} volatility, and mortgage
hedging raise it?
\end{interviewq}

\begin{interviewq}[$\star\star\star$ \iqroles{developer, researcher}]\label{iq:m2:the-rates-options-market:6}
Implied \omterm{def:m2:the-rates-options-market:normal}{normal volatility} must be computed for a million options a second.
How would you do it, and how would you test it?
\end{interviewq}
